\subsection{COMMUTATION FACTORS}

When in this paragraph we speak of graduations without specifying them, we shall always mean graduations of type \(\Delta\) , where \(\Delta\) is a commutative group written additively. In this paragraph, a commutation factor over \(\Delta\) with values in Z is called a commutation factor over \(\Delta\) ( \(\S\) 4, no. 7, Definition 6). A commutation factor over \(\Delta\) is therefore identified with a mapping \(\varepsilon\colon(\alpha,\beta)\mapsto\varepsilon_{\alpha\beta}=\varepsilon(\alpha,\beta)\) of \(\Delta\times\Delta\) into the multiplicative group \(\{-1,1\}\) such that for \(\alpha,\alpha^{\prime},\beta,\beta^{\prime}\) in \(\Delta\) ,

\[
\left\{ \begin{array}{l} \varepsilon (\alpha + \alpha^ {\prime}, \beta) = \varepsilon (\alpha , \beta) \varepsilon (\alpha^ {\prime}, \beta) \\ \varepsilon (\alpha , \beta + \beta^ {\prime}) = \varepsilon (\alpha , \beta) \varepsilon (\alpha , \beta^ {\prime}) \\ \varepsilon (\beta , \alpha) = \varepsilon (\alpha , \beta). \end{array} \right.\tag{1}
\]

It follows that \(\varepsilon(2\alpha, \beta) = \varepsilon(\alpha, 2\beta) = 1\) .

When \(\Delta = \mathbf{Z}\) , every commutation factor \(\varepsilon\) is determined by giving \(\varepsilon(1, 1)\) ; there are therefore only two such factors, the first defined by

\[
\varepsilon (p, q) = 1 \quad \text { for } \quad p, q \text { in } \mathbf {Z}\tag{2}
\]

and the second by

\[
\varepsilon (p, q) = (- 1) ^ {p q} \quad \text { for } \quad p, q \text { in } \mathbf {Z}.\tag{3}
\]

